\documentclass{situationist-preprint}

\title{Confabulation and Self-Report}
\msid{MND-20-b}
\msfamily{MND-20}
\msdiscipline{Philosophy of Mind}
\mstype{Companion manuscript}
\mskeywords{confabulation; introspection; self-report; causal inference; split brain}

\begin{document}
\maketitle

\begin{abstract}
The literature on confabulation is often read as showing that self-report is unreliable across the board. We argue that the confabulations documented in normal subjects, in choice blindness and in split-brain patients have a common structure: they are reports about the causes or reasons of a state or action, produced by causal inference from a theory and presented in the first person. Reports about what one is currently experiencing are not of this kind. Confabulation therefore confirms, rather than threatens, the conclusion of MND-20-a that introspection is reliable about current phenomenal states and unreliable about their causes.
\end{abstract}

\section{Three cases}

In the studies reviewed by \citet{Nisbett1977}, shoppers who chose the rightmost of four identical stockings explained their choice by its quality, not its position; subjects' reports of why they acted matched the predictions of observers given only a description of the situation. In choice blindness, subjects who were covertly handed the photograph they had not chosen gave reasons for preferring it \citep{Johansson2005}. In split-brain patients, the left hemisphere, shown only part of what prompted an action of the left hand, produced a fluent explanation of the action that did not mention its actual cause \citep{Gazzaniga1978}.

\begin{proposition}\label{p:struct}
In each case the confabulated report is a report about the cause of, or reason for, a state or action, and it is formed by inference from a theory of what would explain it.
\end{proposition}

\begin{proof}
The stockings and choice-blindness reports answer the question why one chose; the split-brain reports, why one acted. The match between subjects' and observers' explanations shows that both drew on the same theory of what causes such choices \citep{Nisbett1977}; the split-brain patient's explanation uses only what the left hemisphere could observe of the situation and of the left hand's action. These are the conditions of Principle~P2 of MND-20-a: a causal relation is known through evidence of counterfactual dependence, and in its absence a theory supplies the answer.
\end{proof}

By Lemma~3.2 of MND-20-a, judgments of this kind are no more reliable in the first person than in the third, which is what the evidence shows. None of the cases is a report about the character of a current experience. The shoppers did not misreport how the stockings felt or looked; the split-brain patient's left hemisphere did not misreport what it saw. Such judgments can be formed by embedding the attended experience in the content of the judgment (Principle~P1 of MND-20-a), and Lemma~3.1 there establishes their reliability.

\begin{objection}[Confabulated experience]
Choice-blindness subjects sometimes cite features of the substituted photograph they could not have attended to before, so they misreport their earlier experience.
\end{objection}

\begin{reply}
They misreport an earlier experience, or the reasons for an earlier choice, not a current one. Memory of a past state is not introspection of a current state, and MND-20-a restricts its claim of reliability to current phenomenal states (Definition~2.3); its reply to this case is Objection~1 there.
\end{reply}

\section{Conclusion}

Confabulation is causal theory in the first person. It shows that self-reports about why we feel, choose and act are inferences, and fallible as inferences are; it does not show that we cannot report what we now experience.

\begin{thebibliography}{9}\small
\bibitem[Gazzaniga and LeDoux(1978)]{Gazzaniga1978} Gazzaniga, M.~S. and LeDoux, J.~E. (1978). \emph{The Integrated Mind}. Plenum Press.
\bibitem[Johansson et~al.(2005)]{Johansson2005} Johansson, P., Hall, L., Sikstr\"om, S. and Olsson, A. (2005). Failure to detect mismatches between intention and outcome in a simple decision task. \emph{Science}, 310(5745), 116--119.
\bibitem[Nisbett and Wilson(1977)]{Nisbett1977} Nisbett, R.~E. and Wilson, T.~D. (1977). Telling more than we can know: Verbal reports on mental processes. \emph{Psychological Review}, 84(3), 231--259.
\end{thebibliography}

\end{document}
